\documentclass[10pt,a4paper]{amsart}
\usepackage[margin=19mm]{geometry}
\usepackage{amsmath,amssymb,mathtools}
\usepackage[scr=rsfs,cal=euler]{mathalfa}
\usepackage{xfrac}
\usepackage[unicode,hidelinks,bookmarks=false]{hyperref}
\newcommand{\ie}{i.e.\ }
\newcommand{\cf}{cf.\ }
\newcommand{\resp}{respectively\ }
\newcommand{\Subsection}{\subsection}
\providecommand{\nobreakdash}{\nobreak\hbox{-}}
\setlength{\emergencystretch}{2em}
\multlinegap0pt
\usepackage{xcolor}
\usepackage{amssymb}
\newtheorem{lemma}{Lemma}
\title{Two new proofs of Chui's Conjecture in Weighted Bergman Spaces}
\author{Georgia Corbett}
\date{Source: arXiv:2609.21121v1 (CC BY 4.0)}
\begin{document}
\maketitle

\noindent\emph{Source and licence.} Two new proofs of Chui's Conjecture in Weighted Bergman Spaces. Version arXiv:2609.21121v1 (CC BY 4.0), distributed by arXiv under the Creative Commons Attribution 4.0 International licence, http://creativecommons.org/licenses/by/4.0/, as recorded in the arXiv licence metadata for this exact version and shown on the abstract page https://arxiv.org/abs/2609.21121v1. Full licence notice: \texttt{licenses/arxiv-2609.21121-CC-BY-4.0.txt}.\par\smallskip

\noindent\textit{Editorial scope.} Counted unit: the article's lemma dim (source0.tex lines 1213--1310). The article's complete original proof of it is delivered with it (source0.tex lines 1312--1569, one \texttt{proof} environment of 258 lines). No mathematics is written, repaired or re-derived: the statement and the proof are the article's own lines, in the article's own notation, and no retained line is substituted or re-flowed.  No further source-local context is needed: every cross-reference of every retained line resolves inside the two delivered spans.  Nothing was omitted from any retained span, so the package carries no omission record.  No retained line contains a citation, so the wrapper carries no bibliography and no bibliography entry.  No retained line contains a figure environment, an image-inclusion command or drawing code, so no image is delivered and the package declares no asset dependency.  Editorial formatting only, in addition: an AMS \texttt{amsart} preamble that restates the article's own declarations and macros used by the retained lines, namely the environments \texttt{lemma} on separate counters, together with the standard packages those lines need.  The retained environments print in this document as Lemma 1 in order of appearance, whereas the article prints the same items with the numbering of its own full document.  The complete native source of the article as distributed in the arXiv e-print for this exact version is bundled unchanged as \texttt{sources/210939/source0.tex}.
\begin{lemma}
Let \[A =\left[I - T_N\;\middle|\;T_N^{-1} - T_N\middle|\;\cdots \;\middle|\; T_N^{-(M-2)} - T_N \;\middle|\; T_N^{-(M-1)} \right]. \]
Then we have
\begin{equation}\label{n even dim}
    \dim\ker(A)=MN-N=(M-1)N.
\end{equation}

We claim that
\begin{equation}\label{n even span}
    \ker(A)= \operatorname{span} \{V_0,V_1,\dots,V_{(M-1)N-1}\}.
\end{equation}

The vector $V_0$ is defined as in the odd case:
\[
V_0=
[1\;1\;\dots\;1
\mid
\mathbf{0}
\mid
\dots
\mid
\mathbf{0}]^T.
\]

For
\[
2\leq m\leq M-1,
\qquad
1\leq k\leq N,
\]
define
\[
V_{(m-2)N+k}
=
[
\mathbf{v}_{(m-2)N+k}^{(1)}
\mid
\mathbf{v}_{(m-2)N+k}^{(2)}
\mid
\dots
\mid
\mathbf{v}_{(m-2)N+k}^{(M)}
]^T,
\]
where
\[
\mathbf{v}_{(m-2)N+k}^{(1)}
=-\sum_{r=0}^{m-1} \mathbf{e}_{k+r}
%-\mathbf{e}_k-\mathbf{e}_{k+1}-\dots-\mathbf{e}_{k+(m-1)},
\]
\[
\mathbf{v}_{(m-2)N+k}^{(m)}
=
\mathbf{e}_k,
\]
and
\[
\mathbf{v}_{(m-2)N+k}^{(p)}
=
\mathbf{0}
\qquad
\text{for all }p\neq 1,m.
\]

Lastly, define  for  $1\leq k\leq N-1$
\[
V_{(M-2)N+k}
=
[
\mathbf{v}_{(M-2)N+k}^{(1)}
\mid
\mathbf{v}_{(M-2)N+k}^{(2)}
\mid
\dots
\mid
\mathbf{v}_{(M-2)N+k}^{(M)}
]^T,
\]
where
\[
\mathbf{v}_{(M-2)N+k}^{(1)}
=
\mathbf{e}_{k+M},
\]
\[
\mathbf{v}_{(M-2)N+k}^{(M)}
=
\mathbf{e}_k-\mathbf{e}_{k+1},
\]
and
\[
\mathbf{v}_{(M-2)N+k}^{(p)}
=
\mathbf{0}
\qquad
\text{for all }p\neq 1,m .
\]
\end{lemma}

\begin{proof}
\noindent\textbf{Proof of \eqref{n even dim}:}

Following the structure of the proof of the odd case, we see that
\[
\dim\ker(A)
=
MN-N
=
(M-1)N.
\]

Thus, to characterize $\ker(A)$, we consider the proposed vectors
\[
V_0,V_1,\dots,V_{(M-1)N-1}
\]
and show that they are linearly independent and belong to $\ker(A)$.

\noindent\textbf{Proof of \eqref{n even span}:}

Our goal is to show that
\[
V_i\in\ker(A)
\qquad
\text{for all }
i\in\{0,\dots,(M-1)N-1\}.
\]

From the odd case, we already have
\[
V_0\in\ker(A).
\]

\noindent Now consider $V_{(m-2)N+k}$ for $2\leq m\leq M-1, \quad 1\leq k\leq N.$
From the block structure,
\[
A\mathbf{V}_{(m-2)N+k}
=
(I-T)\mathbf{v}_{(m-2)N+k}^{(1)}
+
(T^{-(m-1)}-T)\mathbf{v}_{(m-2)N+k}^{(m)}.
\]

From the definitions of the vectors
% \[
% \mathbf{v}_{(m-2)N+k}^{(1)}
% =
% -\mathbf{e}_k
% -\mathbf{e}_{k+1}
% -\dots
% -\mathbf{e}_{k+(m-1)}
% \]
% and
% \[
% \mathbf{v}_{(m-2)N+k}^{\,m}
% =
% \mathbf{e}_k,
% \]
we obtain
\begin{align*}
A\mathbf{V}_{(m-2)N+k}
&=
(I-T)\mathbf{v}_{(m-2)N+k}^{(1)}
+
(T^{-(m-1)}-T)\mathbf{v}_{(m-2)N+k}^{(m)}\\
&=
(I-T)
\left(-\sum_{r=0}^{m-1} \mathbf{e}_{k+r}
% -\mathbf{e}_k
% -\mathbf{e}_{k+1}
% -\dots
% -\mathbf{e}_{k+(m-1)}
\right)
 +
(T^{-(m-1)}-T)\mathbf{e}_k.
\end{align*}

Now notice that
\[
T\mathbf{e}_k=\mathbf{e}_{k-1}
\qquad\text{and}\qquad
T^{-p}\mathbf{e}_k=\mathbf{e}_{k+p}.
\]

Thus
\begin{align*}
A\mathbf{V}_{(m-2)N+k}
&=
(I-T)
\left(
-\mathbf{e}_k
-\mathbf{e}_{k+1}
-\dots
-\mathbf{e}_{k+(m-1)}
\right)
+
(T^{-(m-1)}-T)\mathbf{e}_k\\
&=
\left[-\mathbf{e}_k
-\mathbf{e}_{k+1}
-\dots
-\mathbf{e}_{k+(m-2)}
-\mathbf{e}_{k+(m-1)}\right]\\
&\quad
+\left[\mathbf{e}_{k-1}
+\mathbf{e}_k
+\mathbf{e}_{k+1}
+\dots
+\mathbf{e}_{k+(m-2)}\right]
+\mathbf{e}_{k+(m-1)}
-\mathbf{e}_{k-1}\\
&=
0.
\end{align*}

Lastly, consider $V_{(M-2)N+k}$ for
\[
1\leq k\leq N.
\]
Then
\begin{align*}
A\mathbf{V}_{(M-2)N+k}
&=
(I-T)\mathbf{v}_{(M-2)N+k}^{(1)}
+
T^{-(M-1)}
\mathbf{v}_{(M-2)N+k}^{(M)}\\
&=
(I-T)\mathbf{e}_{k+M}
+
T^{-(M-1)}
(\mathbf{e}_k-\mathbf{e}_{k+1})\\
&=
\mathbf{e}_{k+M}
-
\mathbf{e}_{k+M-1}
+
\mathbf{e}_{k+M-1}
-
\mathbf{e}_{k+M}\\
&=
0.
\end{align*}

Therefore, all of the proposed vectors belong to $\ker(A)$.

It remains to show that the proposed kernel vectors are linearly independent.

Consider the linear combination
\[\sum_{i=0}^{(M-1)N-1}a_iV_i\]
with coefficients $a_i\in\mathbb{R}.$

\noindent By construction, in the blocks
\[2\leq m\leq M-1,\]
the only nonzero vectors are
\[
\mathbf{v}_{(m-2)N+k}^{(m)}
=
\mathbf{e}_k.
\]
Therefore,
\[
a_1=a_2=\dots=a_{(M-2)N}=0.
\]

Notice that the remaining coupled entries of the linear combination are in the first entries and the $(M-1)N+1^{th}$ entry through $MN^{th}$ entry.

Since $a_1=a_2=\dots=a_{(M-2)N}=0,$ the first entry then reduces to
\[
a_0-a_{(M-2)N+(M+1)}=0,
\]

\[
\mathbf{v}_{(M-2)N+(M+1)}
=
\mathbf{e}_{(M+1)+M}
=
\mathbf{e}_{N+1}.
\]

Next, we show that the coefficients $a_{(M-1)N+1}$ through $a_{MN}$ have a recursive relationship coming from
\[
\mathbf{v}_{(M-2)N+k}^{(M)}
=
\mathbf{e}_k-\mathbf{e}_{k+1}
\]
forcing the coefficients to be 0.

From the definition of the last block of vectors we first have
\[ a_{(M-2)N+1}=0.\]
Then utilizing the recursive relationship we obtain
\begin{align*}
    a_{(M-2)N+2}-a_{(M-2)N+1}=0 &\quad\Rightarrow \quad a_{(M-2)N+2}=0\\
    a_{(M-2)N+3}-a_{(M-2)N+2}=0\ &\quad\Rightarrow \quad a_{(M-2)N+3}=0\\
    \vdots\\
    a_{(M-2)N+(N-1)}-a_{(M-2)N+(N-2)}=0\ &\quad\Rightarrow \quad a_{(M-2)N+(N-1)}=0\\
\end{align*}

Lastly, we see from the definition of the vectors \[a_{(M-2)N+(N-1)}=0.\]

% Then setting coefficient $a_{(M-2)N+2}$ to be zero:
% \[a_{(M-2)N+2}-a_{(M-2)N+1}=0 \quad\Rightarrow \quad a_{(M-2)N+2}=0.\]

% In the same manner, we argue that all coefficients $a_{(M-2)N+k}$ for $0\leq k\leq N-1$
% \textcolor{red}{
% \begin{align*}
% \textbf{Row } \mathbf{(M-1)N+1}: \quad
% & a_{(M-2)N+1}=0,
% \\[1ex]
% \textbf{Row } \mathbf{(M-1)N+2}: \quad
% & a_{(M-2)N+2}-a_{(M-2)N+1}=0
% \\
% &\qquad\Rightarrow
% a_{(M-2)N+2}=0,
% \\[1ex]
% \textbf{Row } \mathbf{(M-1)N+3}: \quad
% & a_{(M-2)N+3}-a_{(M-2)N+2}=0
% \\
% &\qquad\Rightarrow
% a_{(M-2)N+3}=0,
% \\
% &\quad\vdots
% \\
% \textbf{Row } \mathbf{(M-1)N+N+1}: \quad
% & a_{(M-2)N+(N-1)}
% -
% a_{(M-2)N+(N-2)}
% =
% 0
% \\
% &\qquad\Rightarrow
% a_{(M-2)N+(N-1)}=0.
% \end{align*}
% This then forces
% \[
% \textbf{Row 1:}
% \qquad
% a_0-a_{(M-2)N}
% =
% a_0
% =
% 0.
% \]
Thus,
\[
a_0
=
a_1
=
\dots
=
a_{(M-2)N+(N-1)}
=
0.
\]

Therefore, the proposed vectors are linearly independent.
\end{proof}

\end{document}
